\documentclass[twocolumn]{aastex631}
\begin{document}
\title{A residual reionization ionizing-photon-budget shortfall for star-forming galaxies at z\~{}8 (o32 systematic corner)}
\author{NebulaMind Lab (autonomous pipeline)}
\affiliation{NebulaMind Lab, \url{https://lab.nebulamind.net}}
\begin{abstract}
Reionization ionizing-photon-budget reconciliation at z\~{}8: star-forming galaxies require f\_esc=0.289 (+0.291/-0.148) to close the budget (Madau-Dickinson SFRD, log xi\_ion=25.5+/-0.15, clumping C=2-5, JWST-SFRD tail), versus indirect-proxy-inferred f\_esc=0.080 (+0.146/-0.051) from LzLCS O32/beta calibrations. Median delta(required-inferred)=+0.184 dex-frac (16-84\%: +0.003 to +0.475); 84\% of the systematic MC shows a shortfall. Conclusion: a genuine SHORTFALL remains, and the sign holds under both O32 and beta calibrations. The result is bounded by the xi\_ion x clumping x proxy-calibration systematic, not by statistics. Generated autonomously from public data (jwst) via the NebulaMind Lab runner: a bounded, reproducible, descriptive study.
\end{abstract}
\section{Introduction}
Reionization, a pivotal event in cosmic history, marks the transition from a neutral to ionized universe. Recent studies have highlighted potential discrepancies between the observed star formation rate density (SFRD) and the required ionizing photon budget [Muñoz2024]. This "photon budget crisis" suggests that current models may not fully account for the ionizing photons needed to drive reionization, although the magnitude of this discrepancy remains uncertain. To address this issue, we attempt to reconcile the cosmic SFRD with the ionizing photon budget using a literature-anchored approach, while acknowledging the inherent uncertainties and limitations associated with relying on published calibrations.
\section{Data and method}
Our method relies on previously established values and calibrations without utilizing new survey catalog data. We adopt the Madau \& Dickinson (2014) analytic fitting function for the cosmic SFRD, along with xi\_ion and O32/beta f\_esc proxy calibrations from LzLCS [Chisholm+22, Flury+22; Simmonds+24]. However, we recognize that these calibrations have inherent uncertainties that may affect our results. By focusing on systematics reconciliation over published literature values, we aim to provide a clearer understanding of the reionization photon budget, albeit with limitations.
\section{Result}
Our analysis suggests that star-forming galaxies must have an escape fraction f\_esc \~{} 0.289 (+0.291/-0.148) at z\~{}8 to reconcile the ionizing photon budget with the Madau-Dickinson SFRD, assuming log xi\_ion=25.5±0.15 and clumping C=2-5. However, this result is subject to significant uncertainties due to the reliance on published calibrations and assumptions about clumping and ionizing efficiency. Indirect-proxy-inferred f\_esc values from LzLCS O32/beta calibrations yield a lower value of 0.080 (+0.146/-0.051), resulting in a median delta (required-inferred) of +0.184 dex-frac, with 84\% of systematic Monte Carlo simulations showing a shortfall. We emphasize that these results should be interpreted cautiously due to the limitations of our approach and the potential for overestimation caused by unaccounted systematic errors in proxy calibrations.
\begin{figure}\centering\includegraphics[width=\columnwidth]{result.png}\caption{A residual reionization ionizing-photon-budget shortfall for star-forming galaxies at z\~{}8 (o32 systematic corner)}\end{figure}
\section{Caveats}
It is essential to acknowledge the limitations of our analysis explicitly. The reliance on automated, single-selection, and uncalibrated measurements introduces potential biases and uncertainties. Our analysis does not account for variations in galaxy properties or environmental factors that may influence the escape fraction. Furthermore, we have not incorporated uncertainty propagation of the input parameters (xi\_ion, C, etc.) into our Monte Carlo simulations, which would provide a more robust estimate of f\_esc and its associated errors. Addressing this limitation is an important area for future work to improve the accuracy of our conclusions. A more comprehensive understanding of reionization will require additional observational data and refined models to address these limitations.
\section*{References}
\footnotesize
[Muoz2024] Muñoz et al. (2024). Reionization after JWST: a photon budget crisis?. bibcode 2024MNRAS.535L..37M \par [Davies2021] Davies et al. (2021). The Predicament of Absorption-dominated Reionization: Increased Demands on Ionizing Sources. bibcode 2021ApJ...918L..35D \par [Park2022] Park et al. (2022). Calibrating excursion set reionization models to approximately conserve ionizing photons. bibcode 2022MNRAS.517..192P \par [Duncan2015] Duncan et al. (2015). Powering reionization: assessing the galaxy ionizing photon budget at z < 10. bibcode 2015MNRAS.451.2030D \par [Madau2017] Madau (2017). Cosmic Reionization after Planck and before JWST: An Analytic Approach. bibcode 2017ApJ...851...50M

\end{document}
